\documentclass[12pt,a4paper]{ctexart}

\usepackage[top=2.5cm, bottom=2.5cm, left=2.8cm, right=2.8cm]{geometry}
\usepackage{amsmath,amssymb,amsthm}
\usepackage{booktabs}
\usepackage{enumitem}
\usepackage[colorlinks=true,linkcolor=blue,citecolor=blue]{hyperref}
\usepackage{xcolor}

\newcommand{\tdet}{t_{\text{det}}}
\newcommand{\Teff}{T_{\text{eff}}}
\newcommand{\MA}{|\mathbf{M}\mathbf{A}|}

\title{\textbf{问题一：遮蔽模型与二分求根法}\\[5pt]
       \large 烟幕干扰弹有效遮蔽时长的精确计算}
\author{}
\date{\today}

\begin{document}

\maketitle

\section{问题描述与运动学模型}

\subsection{场景与坐标系}

采用右手直角坐标系：假目标 $\mathbf{O}_{\text{decoy}} = (0,0,0)$ 为原点，$X$ 轴正方向指向导弹来袭方向，$Y$ 轴正方向与 $X$ 轴在水平面内正交，$Z$ 轴正方向竖直向上。真目标是需要保护的圆柱体：底面圆心 $\mathbf{O}_1 = (0, 200, 0)$（位于 $Y$ 轴正半轴，距原点 $200\ \mathrm{m}$），半径 $R_c = 7\ \mathrm{m}$，高度 $H_c = 10\ \mathrm{m}$。

导弹 M1 被雷达发现时位于 $\mathbf{M}_0 = (20000, 0, 2000)$，即 $X$ 轴正半轴方向约 $20\ \mathrm{km}$、高度 $2\ \mathrm{km}$ 处。它以恒定速率 $V_M = 300\ \mathrm{m/s}$ 沿直线直指假目标原点。方向单位向量为：
\begin{equation}
\hat{\mathbf{u}}_m = \frac{\mathbf{O}_{\text{decoy}} - \mathbf{M}_0}{\|\mathbf{O}_{\text{decoy}} - \mathbf{M}_0\|}
= \frac{(-20000, 0, -2000)}{\sqrt{20000^2 + 2000^2}} \approx (-0.995, 0, -0.100)
\end{equation}
$\hat{\mathbf{u}}_m$ 的含义：导弹沿 $X$ 轴负方向偏下方俯冲，水平速度分量约 $298.5\ \mathrm{m/s}$，垂直下沉分量约 $30\ \mathrm{m/s}$。导弹位置的时间函数为匀速直线运动：
\begin{equation}
\boxed{\mathbf{M}(t) = \mathbf{M}_0 + \hat{\mathbf{u}}_m \cdot V_M \cdot t}
\end{equation}
导弹从初始位置飞到原点需时 $T_{\text{total}} = \|\mathbf{M}_0\| / V_M = \sqrt{20000^2 + 2000^2}/300 \approx 67.1\ \mathrm{s}$。

无人机 FY1 位于 $\mathbf{P}_{\text{FY1}} = (17800, 0, 1800)$（在导弹起点沿 $X$ 轴负方向约 $2.2\ \mathrm{km}$、同高度偏低 $200\ \mathrm{m}$）。问题一中 FY1 以固定速度 $V_{\text{drone}} = 120\ \mathrm{m/s}$ 朝假目标方向（方向向量 $\hat{\mathbf{u}}_f = (-1, 0, 0)$，即沿 $X$ 轴负方向）水平直线飞行。在受领任务后 $t_{\text{rel}} = 1.5\ \mathrm{s}$ 时刻投放烟幕弹，延时 $t_{\text{delay}} = 3.6\ \mathrm{s}$ 后起爆。

\subsection{烟幕弹平抛运动与起爆点}

炸弹投放后水平方向继承无人机的速度（惯性），垂直方向仅受重力 $g = 9.8\ \mathrm{m/s^2}$ 作用做自由落体。投放点 $\mathbf{R}_{\text{pt}}$（投放瞬间无人机位置）与起爆点 $\mathbf{A}_{\text{det}}$ 分别为：
\begin{align}
\mathbf{R}_{\text{pt}} &= \mathbf{P}_{\text{FY1}} + V_{\text{drone}} \cdot \hat{\mathbf{u}}_f \cdot t_{\text{rel}} = (17620, 0, 1800) \\
\mathbf{A}_{\text{det}} &= \mathbf{R}_{\text{pt}} + V_{\text{drone}} \cdot \hat{\mathbf{u}}_f \cdot t_{\text{delay}} + (0, 0, -\tfrac{1}{2} g t_{\text{delay}}^2)
\end{align}
代入得 $\mathbf{A}_{\text{det}} = (17188, 0, 1736.5)$。物理含义：无人机沿 $X$ 轴负方向飞行 $1.5\ \mathrm{s}$ 后到达投放点，炸弹继续沿 $X$ 轴负方向水平运动 $3.6\ \mathrm{s}$ 同时下落 $0.5 \times 9.8 \times 3.6^2 \approx 63.5\ \mathrm{m}$，在 $(17188, 0, 1736.5)$ 处起爆。

起爆后烟幕瞬时形成半径 $R = 10\ \mathrm{m}$ 的球体，以 $V_S = 3\ \mathrm{m/s}$ 匀速下沉，有效期 $T_L = 20\ \mathrm{s}$。起爆时刻 $\tdet = t_{\text{rel}} + t_{\text{delay}} = 5.1\ \mathrm{s}$。烟幕球心位置的时间函数为：
\begin{equation}
\boxed{\mathbf{A}(t) = \mathbf{A}_{\text{det}} - (0, 0, V_S) \cdot (t - \tdet), \quad t \in [\tdet,\; \tdet + T_L]}
\end{equation}
$t < \tdet$ 时烟幕尚未起爆，$t > \tdet + T_L$ 时烟幕已失效（浓度不足以遮蔽）。目标是计算该固定投放参数下烟幕对 M1 的有效遮蔽总时长 $\Teff$。

\section{二分求根法}

\subsection{遮蔽的物理条件与临界方程}

导弹对真目标的视线被烟幕球遮蔽，需满足两类条件之一。\textbf{C1 条件}：导弹自身处于烟幕球内部，即导弹与烟幕球心距离小于球半径：
\begin{equation}
\text{C1:}\quad \MA(t) := \|\mathbf{M}(t) - \mathbf{A}(t)\| \leq R = 10\ \mathrm{m}
\end{equation}
\textbf{C2 条件}：导弹在球外（$\MA > R$），但烟幕球的角投影覆盖真目标的角投影，且烟幕位于目标和导弹之间。为此引入两个核心量。

\textbf{烟幕球角半径} $\alpha(t)$——从导弹视角看，烟幕球所张的视角半角：
\begin{equation}
\boxed{\alpha(t) = \arcsin\!\left(\frac{R}{\MA(t)}\right)}
\end{equation}
$\alpha(t)$ 随导弹靠近烟幕球（$\MA$ 减小）而增大——物理上，越近的物体看起来越大。当 $\MA = R$ 时 $\alpha = 90^\circ$（C1 边界），当 $\MA \to \infty$ 时 $\alpha \to 0^\circ$。

\textbf{圆柱最大偏角} $\theta_{\max}(t)$——从导弹指向烟幕球心的方向 $\hat{\mathbf{u}}_A = (\mathbf{A} - \mathbf{M})/\MA$ 到圆柱体可见边界上任意点的方向之间的最大夹角。通过在圆柱体表面选取 215 个特征采样点（两侧母线各 21 点、顶圆 49 点、底面前沿 25 点、前表面内部 45 点），将每个采样点 $\mathbf{p}$ 投影到以导弹为中心的单位球面（方向向量 $(\mathbf{p} - \mathbf{M})/\|\mathbf{p} - \mathbf{M}\|$），与 $\hat{\mathbf{u}}_A$ 求夹角后取最大值：
\begin{equation}
\boxed{\theta_{\max}(t) = \max_{\mathbf{p} \in \mathcal{B}}\; \arccos\!\left(\frac{\mathbf{p} - \mathbf{M}}{\|\mathbf{p} - \mathbf{M}\|} \cdot \hat{\mathbf{u}}_A\right)}
\end{equation}
$\theta_{\max}(t)$ 衡量从烟幕球心方向看出去，真目标圆柱体``占据''了多大的视角。当 $\theta_{\max} \leq \alpha$ 时，圆柱的投影完全被烟幕球的投影圆覆盖，视线被遮挡。

此外还需满足\textbf{几何前置条件}：烟幕球心必须比真目标更靠近导弹，即 $\|\mathbf{M} - \mathbf{O}_1\| > \|\mathbf{A} - \mathbf{O}_1\|$。若烟幕沉到目标后方，即使角度条件满足也无法遮蔽。

因此 $\text{shielded}(t)$ 是分段常数函数——仅在以下三类方程取等号的临界时刻切换：$\MA = R$（C1 边界）、$\theta_{\max} = \alpha$（C2 边界）、$\|\mathbf{M}\mathbf{O}_1\| = \|\mathbf{A}\mathbf{O}_1\|$（几何失效边界）。策略是\textbf{不逐点扫描，而是直接求解这三个方程的精确根}，由根的位置确定遮蔽起止。

\subsection{C1 边界——二次方程解析解}

展开 $\MA^2 = R^2$。将烟幕运动重写为 $\mathbf{A}(t) = \tilde{\mathbf{A}}_0 - (0,0,V_S)\,t$，其中 $\tilde{\mathbf{A}}_0 = \mathbf{A}_{\text{det}} + (0,0,V_S)\tdet$ 是将烟幕的下沉时移吸收到常数项后的等效起始位置。导弹与烟幕的相对运动化为匀速：
\begin{equation}
\mathbf{M}(t) - \mathbf{A}(t) = \underbrace{(\mathbf{M}_0 - \tilde{\mathbf{A}}_0)}_{\mathbf{d}_0} + \underbrace{(\hat{\mathbf{u}}_m V_M + (0,0,V_S))}_{\mathbf{v}_{\text{rel}}} \cdot t
\end{equation}
$\mathbf{d}_0$ 是起爆瞬间弹-目相对位置（$\tdet$ 时刻 $\mathbf{M} - \mathbf{A}$ 的值），$\mathbf{v}_{\text{rel}}$ 是相对速度（导弹水平俯冲速度 $\hat{\mathbf{u}}_m V_M$ 与烟幕下沉速度 $(0,0,V_S)$ 的矢量和）。注意 $V_S = 3 \ll V_M = 300$，故 $\mathbf{v}_{\text{rel}} \approx \hat{\mathbf{u}}_m V_M$，相对运动主要由导弹驱动。$\|\mathbf{d}_0 + \mathbf{v}_{\text{rel}} t\|^2 = R^2$ 展开为标准二次方程：
\begin{equation}
\boxed{a_2 t^2 + a_1 t + a_0 = 0}
\end{equation}
其中 $a_2 = \|\mathbf{v}_{\text{rel}}\|^2$（相对速率平方，恒正，保证二次函数开口向上），$a_1 = 2(\mathbf{d}_0 \cdot \mathbf{v}_{\text{rel}})$（初始相对位置在相对速度方向的投影——若相对速度指向初始位置则 $a_1 < 0$，导弹靠近烟幕），$a_0 = \|\mathbf{d}_0\|^2 - R^2$（初始距离平方与球半径平方之差——若起爆时导弹已在球内则 $a_0 < 0$）。判别式 $\Delta = a_1^2 - 4a_2a_0$，当 $\Delta \geq 0$ 时两根：
\begin{equation}
t_{\text{enter}},\; t_{\text{exit}} = \frac{-a_1 \mp \sqrt{\Delta}}{2a_2}
\end{equation}
仅保留在有效区间 $(\tdet,\; \tdet + T_L)$ 内的根，分别对应导弹\textbf{进入}烟幕球和\textbf{离开}烟幕球的精确时刻。此方程为完全解析解，精度仅受浮点舍入误差（$\sim 10^{-15}$）限制。

\subsection{C2 边界——二分求根}

C2 边界方程 $\theta_{\max}(t) = \alpha(t)$ 定义函数：
\begin{equation}
\boxed{f_{\text{C2}}(t) = \alpha(t) - \theta_{\max}(t)}
\end{equation}
$f_{\text{C2}}(t) > 0$ 时 $\alpha$ 大于 $\theta_{\max}$，即烟幕投影覆盖圆柱投影，角度遮蔽成立。$f_{\text{C2}}$ 是 $t$ 的连续函数（$\alpha$ 和 $\theta_{\max}$ 分别由 $\arcsin$ 和 $\max$ 定义，均为连续映射）。$\theta_{\max}$ 含 $\max$ 操作而无封闭解析式，故须用数值求根。算法分两步：

\textbf{第一步——粗扫描}：在 $[\tdet, \tdet + T_L]$ 区间上生成 40 个等距时间点 $\{t_k\}_{k=0}^{39}$（间距 $20/40 = 0.5\ \mathrm{s}$），逐一计算 $f_{\text{C2}}(t_k)$。每次计算需运行 $\theta_{\max}$ 的 215 点投影，40 次共计 $40 \times 215 = 8600$ 次向量运算。\textbf{第二步——精确定位}：检测相邻点间的符号变化——若 $f_{\text{C2}}(t_k) \cdot f_{\text{C2}}(t_{k+1}) < 0$，说明该小区间内存在奇数个根。对每个异号区间应用二分法 $\operatorname{bisect}(f_{\text{C2}}, t_k, t_{k+1}, \text{xtol}=10^{-6})$。$f_{\text{C2}}$ 在单根邻域内单调（物理上方向关系单调变化），二分法 2--3 次迭代内收敛至 $10^{-6}\ \mathrm{s}$（$1\ \mu\mathrm{s}$ 量级）。40 点间距 0.5s 的扫描密度对至多两次状态切换（起和止）的 20s 窗口足够——不会遗漏区间的根。

\subsection{几何失效边界——二次方程解析解}

烟幕球心与真目标等距时，烟幕不再位于导弹和目标之间。方程 $\|\mathbf{M} - \mathbf{O}_1\|^2 = \|\mathbf{A} - \mathbf{O}_1\|^2$，记导弹速度向量 $\mathbf{v}_M = \hat{\mathbf{u}}_m V_M$，两侧展开：
\begin{align}
\|\mathbf{M} - \mathbf{O}_1\|^2 &= \|\mathbf{M}_0 - \mathbf{O}_1\|^2 + 2(\mathbf{M}_0 - \mathbf{O}_1)\!\cdot\!\mathbf{v}_M \cdot t + \|\mathbf{v}_M\|^2 t^2 \\
\|\mathbf{A} - \mathbf{O}_1\|^2 &= \|\tilde{\mathbf{A}}_0 - \mathbf{O}_1\|^2 - 2(\tilde{\mathbf{A}}_0 - \mathbf{O}_1)\!\cdot\!(0,0,V_S) \cdot t + V_S^2 t^2
\end{align}
令两者相等（消去 $t^2$ 项时注意两侧 $t^2$ 系数分别为 $\|\mathbf{v}_M\|^2$ 和 $V_S^2$，不相等），整理为标准二次型：
\begin{equation}
\boxed{a_2^g t^2 + a_1^g t + a_0^g = 0}
\end{equation}
其中 $a_2^g = \|\mathbf{v}_M\|^2 - V_S^2$（$300^2 - 3^2 \gg 0$，开口向上），$a_1^g = 2[(\mathbf{M}_0 - \mathbf{O}_1)\cdot\mathbf{v}_M + (\tilde{\mathbf{A}}_0 - \mathbf{O}_1)\cdot(0,0,V_S)]$，$a_0^g = \|\mathbf{M}_0 - \mathbf{O}_1\|^2 - \|\tilde{\mathbf{A}}_0 - \mathbf{O}_1\|^2$（起爆瞬间导弹与烟幕到目标距离平方之差）。判别式 $\Delta_g = (a_1^g)^2 - 4a_2^g a_0^g$。根在有效区间 $(\tdet, \tdet+T_L)$ 内且满足 $\|\mathbf{M} - \mathbf{O}_1\| \leq \|\mathbf{A} - \mathbf{O}_1\|$ 的一侧（即烟幕沉至目标后方之后），取最早者为几何失效时刻 $t_{\text{geom}}$。

\section{遮蔽时段构造与数值结果}

\subsection{区间并集与总时长}

将三个方程求出的所有根按时间升序排列后，遮蔽时段由以下规则构造。若 C2 有根（说明存在角度遮蔽阶段），则以第一个 C2 根 $t_{\text{C2\_start}}$ 为起点，以 C1 入球时刻 $t_{\text{C1\_enter}}$ 为终点（若无 C1 根则以 $t_{\text{geom}}$ 或 $\tdet+T_L$ 为终点）。若 C1 有根（说明导弹进入烟幕球），则以入球时刻为起点，以 $t_{\text{end}} = \min(t_{\text{C1\_exit}},\; t_{\text{geom}},\; \tdet+T_L)$ 为终点——即球内遮蔽在导弹出球、烟幕沉至目标后方、或烟幕失效三者中最早发生的时刻结束。总遮蔽时长为各不重叠时段长度之和（并集）：
\begin{equation}
\boxed{\Teff = (t_{\text{C1\_enter}} - t_{\text{C2\_start}}) + (t_{\text{end}} - t_{\text{C1\_enter}})}
\end{equation}

\textbf{退化处理}：若三类方程均无根落在有效区间内（参数偏离可行域时可能发生），则回退至以 0.1s 步长在 $[\tdet, \tdet+T_L]$ 上离散扫描 $\text{shielded}(t)$——每步判定导弹是否在球内、是否满足 $\theta_{\max} \leq \alpha$ 且 $|\mathbf{M}\mathbf{O}_1| > |\mathbf{A}\mathbf{O}_1|$，累加遮蔽步数。这种情况通常对应 $\Teff \approx 0$，优化算法自然避开。

\subsection{问题一的数值结果}

代入问题一的固定参数进行计算。各临界时刻的求根结果：
\begin{center}
\begin{tabular}{c|c|c|c|c}
\toprule
$\tdet$ & $t_{\text{C2\_start}}$ & $t_{\text{C1\_enter}}$ & $t_{\text{geom}}$ & $t_{\text{C1\_exit}}$ \\
\midrule
$5.100$ & $8.057$ & $9.389$ & $9.419$ & $ 9.449$ \\
\bottomrule
\end{tabular}
\end{center}

$t_{\text{end}} = 9.449\ \mathrm{s}$——遮蔽因烟幕沉至目标后方而终止，早于导弹出球时刻，因此导弹出球时遮蔽结束。

\textbf{各阶段时长与物理过程}：
\begin{itemize}
    \item \textbf{起爆至 C2 开始}（$5.100 \to 8.057\ \mathrm{s}$，$2.957\ \mathrm{s}$）：烟幕已起爆但导弹尚远，$\alpha$ 过小不足以覆盖 $\theta_{\max}$，无遮蔽。
    \item \textbf{C2 角度遮蔽阶段}（$8.057 \to 9.389\ \mathrm{s}$，$\mathbf{1.332\ \mathrm{s}}$）：导弹逐渐靠近烟幕球，$\alpha$ 增大至覆盖 $\theta_{\max}$，视线被遮挡。导弹在球外（$\MA > 10\ \mathrm{m}$）。
    \item \textbf{C1 球内遮蔽阶段}（$9.389 \to  9.449\ \mathrm{s}$，$\mathbf{0.060\ \mathrm{s}}$）：导弹以 $300\ \mathrm{m/s}$ 沿 $X$ 轴负方向穿入 $R = 10\ \mathrm{m}$ 的烟幕球。穿越时间在量级上确为 $2R / V_M \approx 0.067\ \mathrm{s}$——弹速极高而球体极小，球内停留极短。实际上导弹并非从球心穿过而仅擦过球体边缘，故穿越时间更短。随后烟幕沉至目标后方，遮蔽终止。
\end{itemize}

\textbf{最终有效遮蔽时长}：由于C2条件失效后C1条件才失效，因此有效时间直接等于$t_{\text{C1\_exit}}$ - $t_{\text{C2\_start}}$
\begin{equation}
\boxed{\Teff = 9.449-8.057 = 1.392\ \mathrm{s}}
\end{equation}
	占导弹全程飞行时间 $67.1\ \mathrm{s}$ 的约 $2.0\%$——在给定的固定飞行参数（无人机 $120\ \mathrm{m/s}$ 直飞假目标、固定投放和起爆延时）下，遮蔽效果极为有限，体现了后续问题中优化参数的必要性。

\end{document}
